\documentclass[11pt]{article}
\usepackage[T1]{fontenc}
\usepackage[utf8]{inputenc}
\usepackage{lmodern,amsmath,amssymb,amsthm,mathtools}
\usepackage[margin=28mm]{geometry}
\usepackage{microtype,xcolor}
\usepackage[authoryear,round]{natbib}
\usepackage{xurl}
\usepackage[colorlinks=true,allcolors=blue!50!black]{hyperref}
\newtheorem{theorem}{Theorem}[section]
\newtheorem{proposition}[theorem]{Proposition}
\newtheorem{lemma}[theorem]{Lemma}
\newtheorem{corollary}[theorem]{Corollary}
\theoremstyle{remark}\newtheorem{remark}[theorem]{Remark}
\newcommand{\OO}{\mathcal O}
\newcommand{\FF}{\mathbf F}
\newcommand{\PP}{\mathbf P}
\newcommand{\ZZ}{\mathbf Z}
\newcommand{\et}{\mathrm{et}}
\newcommand{\GL}{\operatorname{GL}}
\newcommand{\Pic}{\operatorname{Pic}}
\newcommand{\Spec}{\operatorname{Spec}}
\newcommand{\rk}{\operatorname{rk}}
\newcommand{\im}{\operatorname{im}}
\newcommand{\coker}{\operatorname{coker}}
\newcommand{\Sym}{\operatorname{Sym}}
\newcommand{\Tr}{\operatorname{Tr}}
\newcommand{\workbenchurl}{https://github.com/KokunoYumeto/schemes-and-moduli-workbench}
\newcommand{\recordurl}{https://github.com/KokunoYumeto/mathematics-workbenches}
\title{Primitive monodromy in Grothendieck's crystalline programme, after Cartier and Sun}
\author{}
\date{23 November 2026}
\begin{document}
\maketitle
\begin{abstract}
Grothendieck's crystalline notes ask whether the monodromy of a varying hypersurface's surviving Frobenius space is the whole linear group. We determine two primitive-system cases in characteristic two. Every fixed nonordinary abelian surface with an ample globally generated polarization of type $(1,3)$ has relative geometric image $\GL_2(\FF_2)$; its nonzero Lang cover is connected of degree three. On the supersingular product $E^2$, where $E:y^2+y=x^3$, the lines $L_m=\OO(mV+H+\Gamma)$ have type $(1,2m+1)$. For every $m\ge1$, a general smooth member has genus $2m+2$, two-rank $2m$, and exact relative image $S_{2m+1}$ in its permutation quotient representation. Its nonzero Lang cover has exactly $m$ connected components, of degrees $\binom{2m+1}{j}$ for $1\le j\le m$. Thus unrestricted full linear image and nonzero-cover connectedness fail in every odd type at least five in this explicit family. These lines have degree two on an elliptic fibre and are not very ample, so the examples do not refute eventual assertions for powers of a fixed embedding line. The proofs use Cartier multiplication, two explicit Lang restrictions, and a complete elliptic branch series obtained by relative Frobenius. The mathematics and text were produced by GPT-6 Astra Pro and checked within the scopes stated at the end; no human has checked the mathematics.
\end{abstract}

\section{The question, the systems, and the main results}\label{sec:intro}

We compare sharp Frobenius ranks with full monodromy images. Grothendieck asks whether the image ``est-elle \'egale \`a $\operatorname{Aut}(V)$ tout entier?'' in his crystalline-cohomology notes, cote~4, \S8.6, archive sheet~11, PDF image~12 \citep{Grothendieck}. His parameter space consists of hypersurfaces of a fixed projective embedding. Primitive polarizations on an abelian surface provide a precise smaller setting in which to examine the proposed mechanism. A large stable rank alone does not force the nonzero vectors to form one monodromy orbit.

Throughout, $k$ is algebraically closed of characteristic two. An abelian surface $A/k$ is \emph{nonordinary} if its two-rank
\[
 f(A)=\dim_{\FF_2}H^1_{\et}(A,\FF_2)
\]
is less than two. For a smooth proper connected curve $C$, write $f(C)=\dim_{\FF_2}H^1_{\et}(C,\FF_2)$. If $L$ is ample, set $W=H^0(A,L)$ and $X=\PP(W)$, with points the one-dimensional spaces of defining sections. We say that $L$ has type $(1,n)$, for odd $n$, when its polarization has elementary divisors $1,n$; its kernel is then an \'etale group of type $(\ZZ/n)^2$, $L^2=2n$, and $h^0(L)=n$.

For a smooth member $C\in|L|$, the relative coherent space is
\[
 Q_C=H^1(C,\OO_C)/H^1(A,\OO_A).
\]
It has dimension $h^0(L)-1$ and inherits a Frobenius-semilinear map. The \emph{relative Lang system} is the $\FF_2$-local system of Frobenius-fixed vectors in its stable part on an open of constant stable rank. On the locus where Frobenius is bijective on $Q_C$, it is simply $\ker(F-1)$ on that bundle. Its \emph{nonzero Lang cover} is the finite \'etale cover of nonzero vectors. Geometric monodromy means the image of the geometric fundamental group on one fibre; connected components of this cover correspond to its orbits on nonzero vectors.

\begin{theorem}[Every fixed globally generated type $(1,3)$ system]\label{thm:three}
Let $A/k$ be a fixed nonordinary abelian surface, and let $L$ be ample, globally generated, and of type $(1,3)$. The open $U\subset|L|$ of smooth members with bijective relative Frobenius is nonempty. Its rank-two relative Lang system has geometric monodromy
\[
                         \GL_2(\FF_2)=S_3.
\]
The nonzero Lang cover is connected of degree three. A general smooth member has genus four and two-rank $f(A)+2$.
\end{theorem}

Thus the image is full on every fixed surface satisfying the assumptions. The theorem does not replace that quantifier by a generic point of a moduli space. Global generation is an explicit hypothesis.

For the second result, let $E$ be the smooth projective elliptic curve
\[
                         y^2+y=x^3,
\]
with origin $O$ at infinity. On $A=E\times E$ put
\begin{equation}\label{eq:lines}
 \begin{split}
 V&=\{O\}\times E,\qquad H=E\times\{O\},\qquad
 \Gamma=\{(a,b):a+b=O\},\\
 L_m&=\OO_A(mV+H+\Gamma),\qquad n=2m+1.
 \end{split}
\end{equation}
Here $a+b$ uses the elliptic group law, not addition of affine coordinates.

\begin{theorem}[Exact symmetric-group images in every odd type]\label{thm:odd}
For every integer $m\ge1$, the line $L_m$ is ample and globally generated of type $(1,n)$, where $n=2m+1$. It satisfies
\[
 L_m^2=2n,\qquad h^0(L_m)=n,\qquad L_m\cdot V=2,
 \qquad f(A)=0.
\]
The smooth maximal-relative-rank open $U_m\subset|L_m|$ is nonempty. Its curves have genus $n+1$ and two-rank $n-1$, and its geometric relative monodromy image is exactly
\begin{equation}\label{eq:Sn}
 S_n\subset\GL\bigl(\FF_2^n/\FF_2\mathbf1\bigr)
          =\GL_{n-1}(\FF_2),
 \qquad \mathbf1=(1,\ldots,1).
\end{equation}
The nonzero Lang cover has exactly $(n-1)/2$ connected components. Their degrees are
\begin{equation}\label{eq:degrees}
 \binom n1,\binom n2,\ldots,\binom n{(n-1)/2}.
\end{equation}
For $n=3$ the image is $S_3=\GL_2(\FF_2)$. For every odd $n\ge5$ it is a proper subgroup of the full linear group and the nonzero cover is disconnected.
\end{theorem}

For example, type $(1,5)$ has genus six, relative rank four, and image of order $120$, whereas $|\GL_4(\FF_2)|=20160$. Its nonzero cover splits into connected covers of degrees five and ten. Type $(1,9)$, which does not require a prime polarization degree, gives connected components of degrees $9,36,84,126$.

\paragraph{Credit and neighboring work.}
The local Cartier operator and its product formula are classical \citep[0, (2.1.4)--(2.1.13), pp.~514--515]{Illusie}. The Frobenius filtration used to establish the type-$(1,3)$ tensor is the surface case of \citet[Definition~3.1, Lemmas~3.2--3.5 and Theorem~3.6, pp.~5--9]{Sun}; we give the complete surface stability proof. Theta descent, its representation theory, and abelian Riemann--Roch are used at the exact locations of \citet[\S\S8--9]{EGM} specified below. The branch-module calculation is a direct Artin--Schreier argument, stronger than its numerical dimension statement. The Frobenius factorization of a curve linear series is the classical mechanism in \citet[Lemma~53.13.8, full proof]{StacksFrob}; completeness of the descended series is proved here by intersection theory.

For complex type-$(1,3)$ surfaces, \citet[Introduction, Theorems~0.1--0.2, pp.~1--2]{CataneseSernesi} study the singular-divisor discriminant and its Hesse geometry. Our Fermat cubic records nonbijectivity of relative Frobenius in characteristic two. Recent AI-assisted work of \citet[Introduction, Theorems~A--C and declarations]{Zou} studies quasi-$F$-splitting and crystalline slopes of primitive symplectic varieties. It is neighboring Frobenius context and supplies no theorem used here. Bounded searches did not identify these exact fixed-system image statements as one earlier result; no priority claim is made.

\paragraph{Scope and mechanism.}
Theorem~\ref{thm:three} first identifies the actual relative cohomology with an Euler bundle. A cyclic cubic restriction and a quadratic Artin--Schreier restriction force its image to be $S_3$. Theorem~\ref{thm:odd} uses a different mechanism. Every curve in the explicit system maps with degree two to a supersingular elliptic curve. Its \'etale cohomology is the permutation quotient on its branch points. The branch series becomes a complete degree-$n$ series on an elliptic curve after relative Frobenius; ordered divisors in that series form an integral variety, giving the exact group $S_n$.

The lines in \eqref{eq:lines} are not very ample, because their restrictions to $V$ have degree two. They vary with $m$ and retain this fibre degree. Theorem~\ref{thm:odd} therefore disproves unrestricted full-image and connectedness assertions for globally generated primitive systems, but does not refute a sufficiently-large-degree assertion for powers of a fixed embedding line. It makes no claim about arbitrary primitive polarizations on other surfaces.

\section{Relative cohomology, fixed vectors, and smooth members}\label{sec:basic}

We establish the cohomological and geometric inputs shared by the two image calculations. On an abelian surface, an ample line $L$ satisfies $H^i(A,L)=0$ for $i>0$, $\chi(L)=L^2/2$, and $H^1(A,L^{-1})=0$ by Serre duality; see \citet[Theorems~9.11 and 9.14, with their proofs, pp.~131--133]{EGM}.

\begin{lemma}[The relative Euler bundle]\label{lem:euler-bundle}
Let $L$ be ample, $N=h^0(L)$, and $X=\PP(H^0(L))$. For the universal divisor $q:\mathcal C\to X$ there is a Frobenius-compatible exact sequence
\begin{equation}\label{eq:coherent}
 0\longrightarrow H^1(A,\OO_A)\otimes\OO_X
 \longrightarrow R^1q_*\OO_{\mathcal C}
 \longrightarrow\Omega_X^1\longrightarrow0.
\end{equation}
Every effective member is connected. A smooth member has genus $N+1$. At a geometric point where the last Frobenius is bijective,
\[
                         f(C)=f(A)+N-1.
\]
\end{lemma}
\begin{proof}
The universal ideal on $A\times X$ is $L^{-1}\boxtimes\OO_X(-1)$. Push forward its divisor sequence. Vanishing gives the first map of \eqref{eq:coherent} and identifies its quotient with
\[
 \ker\bigl(H^2(A,L^{-1})\otimes\OO_X(-1)
                   \longrightarrow H^2(A,\OO_A)\otimes\OO_X\bigr).
\]
Choose an invariant volume form to identify $H^2(A,\OO_A)$ with $k$. Serre duality makes the displayed map the evaluation row $s:H^0(L)^*\otimes\OO_X(-1)\to\OO_X$. It is surjective at every nonzero section, so its kernel is $\Omega_X^1$. This also gives local freeness and base change. Frobenius respects the divisor sequence and the ambient subbundle.

For an individual effective divisor the degree-zero and degree-one terms give $H^0(C,\OO_C)=k$. Adjunction gives $p_a(C)=1+L^2/2=N+1$.

For a finite-dimensional Frobenius-semilinear space over $k$, the map $F-1$ is surjective \citep[Lemma~59.63.2, full proof]{StacksAS}. One elementary proof splits off the eventual kernel and image. On the nilpotent part, $1-F$ has a finite geometric-series inverse. On the bijective part, in a basis its equations reduce to $x_i^2-\sum b_{ij}x_j=a_i$, with $(b_{ij})$ invertible. The quotient algebra has basis the square-free monomials, and its Jacobian is invertible; every fibre has $2^r$ points. This proves surjectivity and identifies the fixed space with $\FF_2^r$. Applying the snake lemma to $F-1$ on an exact sequence shows that fixed vectors are exact. Artin--Schreier identifies these fixed vectors in coherent $H^1$ with $H^1_{\et}(-,\FF_2)$. Taking fixed vectors in \eqref{eq:coherent} proves the formula.

The same equations on a bundle with invertible linearized Frobenius define a finite \'etale scheme: they are finite free with the indicated monomial basis and invertible Jacobian. Their addition gives a locally constant $\FF_2$-vector space. Artin--Schreier and the cohomology sequence identify this family with the relative curve system. Thus subsequent matrix computations concern the actual relative local system.
\end{proof}

\begin{lemma}[Smooth members for odd primitive type]\label{lem:smooth}
If $L$ is ample and globally generated of type $(1,n)$ with $n$ odd, then $|L|$ has a nonempty open of smooth connected members.
\end{lemma}
\begin{proof}
Choose a cyclic maximal isotropic subgroup $K\subset\ker\lambda_L$ of order $n$. Its theta extension splits: a lift of a generator has $n$th power in $k^*$, and rescaling by an $n$th root makes that power one. This is a genuine linearization, so $L$ descends along the \'etale isogeny $A\to B=A/K$ to an ample line $M$. The exact descent theorem is \citet[Theorem~8.10 and Corollary~8.11, pp.~116--117]{EGM}. We have $\chi(M)=\chi(L)/n=1$, hence $M$ is principal.

Let $\Theta$ be its unique effective divisor. Its square is two. Effective curves on an abelian surface are nef: a translate avoiding any chosen curve computes a nonnegative intersection. An abelian variety has no rational curves, and a map from an elliptic curve to it is a homomorphism after translation. If $\Theta$ has one irreducible component, write $\Theta=aD$. Adjunction makes $D^2$ even; $a^2D^2=2$ forces $a=1$. The arithmetic genus is two. Singularity would make its normalization have genus at most one, whose image would be a smooth elliptic translate, a contradiction. Thus $\Theta$ is smooth of genus two.

If $\Theta=\sum a_iD_i$ has at least two components, positivity of $\Theta\cdot D_i$ and
\[
                         2=\sum a_i(\Theta\cdot D_i)
\]
force precisely two components, each of multiplicity one and of $\Theta$-degree one. Their nonnegative even self-intersections are zero and their mutual intersection is one. Adjunction and the absence of rational curves make them smooth elliptic translates. They meet transversely at one point. Therefore $\Theta$ is reduced nodal. Its \'etale inverse image is a smooth or reduced nodal member $C_0=(s_0=0)$ of $|L|$.

If there are nodes, global generation permits choosing $s_1$ nonzero at every one of them. In the pencil $s_0+zs_1$, a completed node has equation $xy+z u(x,y,z)=0$ with $u$ a unit. Its relative singular equations give $x,y\in(z)$; substitution then gives $z$ times a unit. No component of this singular scheme dominates the parameter trait. Away from the nodes the central member is already smooth. Properness excludes a generic singular point elsewhere. The generic member is therefore smooth, and Lemma~\ref{lem:euler-bundle} makes it connected. This proof uses no unqualified positive-characteristic Bertini assertion.
\end{proof}

\section{The type-\texorpdfstring{$(1,3)$}{(1,3)} Cartier tensor}\label{sec:cartier}

We show why every fixed nonordinary surface in Theorem~\ref{thm:three} gives the same nonzero alternating tensor. The nonvanishing step requires a geometric argument; theta irreducibility alone would not prove it.

\begin{lemma}[Odd Euler characteristic and Frobenius stability]\label{lem:stable}
For a line $L$ on an abelian surface with odd $\chi(L)$, the rank-four bundle $F_*L$ is slope stable for every ample polarization on $A^{(2)}$, where $F:A\to A^{(2)}$ is relative Frobenius.
\end{lemma}
\begin{proof}
The surface case of Sun's canonical filtration has graded terms $L$, $L\otimes\Omega_A^1$, and $L\otimes\bigwedge^2\Omega_A^1$ \citep[Definition~3.1, Lemmas~3.2--3.5 and Theorem~3.6, pp.~5--9]{Sun}. Here it also has a direct construction. The diagonal ideal on $A\times_{A^{(2)}}A$ has local generators $\epsilon_1,\epsilon_2$ with $\epsilon_1^2=\epsilon_2^2=0$. Its graded degrees have ranks $1,2,1$ and transform as $\OO_A,\Omega_A^1,\bigwedge^2\Omega_A^1$. Tensoring with the line on the second factor gives the asserted filtration. Since $\Omega_A^1$ is trivial,
\[
                         T=F^*F_*L\otimes L^{-1}
\]
has a filtration by four copies of $\OO_A$. It is semistable of slope zero for every polarization. A larger-slope subsheaf of $F_*L$ would pull back to a larger-slope subsheaf of $T$, proving semistability of $F_*L$.

Suppose a saturated proper subsheaf $Q\subset F_*L$ of rank $q\in\{1,2,3\}$ has equal slope. Put $S=F^*Q\otimes L^{-1}$ and $R=T/S$. Flatness of Frobenius makes $R$ torsion free, and $S,R$ have slope zero. Intersect $S$ with the trivial-quotient filtration. Each nonzero graded piece is a rank-one subsheaf $I_Z(-D)\subset\OO_A$, with $D$ effective and $Z$ zero dimensional. Its degree is nonpositive. The total degree is zero, so every $D$ is zero. Thus $c_1(S)=0$ and $\chi(S)\le0$.

The inclusion $R^*\subset T^*$ gives the same conclusion $\chi(R^*)\le0$. On the smooth surface, $R^{**}$ is locally free and $R^*=(R^{**})^*$. Serre duality with trivial canonical line gives
\[
 \chi(R)=\chi(R^*)-\operatorname{length}(R^{**}/R)\le0.
\]
Since $\chi(T)=0$, additivity forces $\chi(S)=0$. Write $l=c_1(L)$. The determinant equality gives $F^*c_1(Q)=ql$. Riemann--Roch and $\deg F=4$ yield
\[
 0=\chi(S)=4\chi(Q)-F^*c_1(Q)\cdot l+q\frac{l^2}{2}
          =4\chi(Q)-q\chi(L).
\]
Oddness of $\chi(L)$ forces $4\mid q$, a contradiction. This includes possible zero-dimensional defects of both sheaves.
\end{proof}

\begin{lemma}[Three sections and an intersection bound]\label{lem:divisor}
Let $H$ be ample and globally generated of type $(1,3)$ on an abelian surface. If $h^0(D)\ge3$ for a line $D$, then $D\cdot H\ge6$.
\end{lemma}
\begin{proof}
Choose an effective divisor for $D$. It is nef. If $D^2>0$, Hodge index excludes a curve orthogonal to it, since that curve would have negative square. The surface ampleness criterion makes $D$ ample. Vanishing and Riemann--Roch give $D^2=2h^0(D)\ge6$, and Hodge index with $H^2=6$ gives $D\cdot H\ge6$.

If $D^2=0$, write $D=\sum a_iC_i$. Nefness forces every $C_i^2$ and every $C_i\cdot C_j$ to vanish. Adjunction gives arithmetic genus one. The normalization cannot be rational; hence every $C_i$ is a smooth elliptic translate. Different elliptic subgroup directions have positive intersection, since their addition map is an isogeny onto the surface. All components are therefore fibres of one elliptic quotient $\pi:A\to E'$. We have $D=\pi^*D_0$ with $D_0$ effective. Connected fibres give $h^0(D)=h^0(D_0)=\deg D_0\ge3$. The restriction of $H$ to an elliptic fibre is globally generated and of positive degree, hence has degree at least two. Therefore $D\cdot H\ge2\deg D_0\ge6$.
\end{proof}

\begin{proposition}[The nonzero alternating tensor]\label{prop:tensor}
Under Theorem~\ref{thm:three}'s assumptions, relative Cartier multiplication gives an isomorphism
\[
 \beta:\bigwedge^2 W\longrightarrow H^0(A^{(2)},L^{(2)}),
 \qquad W=H^0(A,L).
\]
There is a level-three basis $e_0,e_1,e_2$ and $c\in k^*$ such that
\begin{equation}\label{eq:tensor}
 \beta(e_i,e_j)=c e_k^{(2)}\quad
                    \text{when }\{i,j,k\}=\{0,1,2\}.
\end{equation}
\end{proposition}
\begin{proof}
Choose invariant volume forms. Finite-flat duality and Cartier trace give a perfect pairing
\[
                 F_*L\otimes F_*L\longrightarrow L^{(2)}.
\]
Locally its scalar part is the coefficient of the top monomial in the basis $1,x,y,xy$ over the Frobenius subring. This gives its perfectness and the projection formula. The invariant top differential has zero Cartier image when $A$ is nonordinary: the top Hasse scalar is the determinant of Frobenius on $H^1(A,\OO_A)$. Hence the pairing vanishes on $(v,v)$ and is alternating. The local Cartier formulas used here are those of \citet[0, (2.1.4)--(2.1.13), pp.~514--515]{Illusie}.

The theta group has perfect commutator pairing on $K(L)\simeq(\ZZ/3)^2$ \citep[Proposition~8.19, Corollary~8.20, Lemma~8.24 and Theorem~8.28, pp.~119--123]{EGM}. Choose lifts of generators with orders three and primitive commutator $\omega$. A clock eigenvector and its shift translates give a basis with shift $e_i\mapsto e_{i+1}$ and clock $e_i\mapsto\omega^i e_i$. The three distinct clock eigenvalues and the transitive shift prove irreducibility of $W$. They also prove irreducibility of $\bigwedge^2W$, of central weight two. The target is the Frobenius twist of $W$ and has the same central weight. Cartier commutes with translation and theta line isomorphisms, so $\beta$ is equivariant, and is either zero or an isomorphism. The clock weights $i+j=2k$ modulo three and the shift force every nonzero map to have \eqref{eq:tensor}.

Suppose $\beta=0$. Let $Q$ be the saturation of the global evaluation image in $F_*L$. Its generic fibre is isotropic in a four-dimensional symplectic space, so its rank $q$ is one or two. It is theta invariant. At rank one, its determinant has the three independent sections of $W$. At rank two, the nonzero equivariant wedge-evaluation map $\bigwedge^2W\to H^0(\det Q)$ is injective by irreducibility. Thus $h^0(\det Q)\ge3$ in both cases. Determinant sections extend across the finite singular set on the smooth surface.

Put $H=L^{(2)}$. The Frobenius filtration, or Riemann--Roch, gives $c_1(F_*L)=2c_1(H)$ numerically. Since $H^2=6$, the slope is three. Lemma~\ref{lem:stable} gives
\[
                         c_1(Q)\cdot H<3q\le6,
\]
contradicting Lemma~\ref{lem:divisor}. Hence $\beta$ is nonzero.
\end{proof}

\section{The Euler matrix and the full rank-two image}\label{sec:three-proof}

We now pass from the geometric tensor to an explicit Lang system and exhibit both needed permutation orders within that same system.

\begin{proposition}[The actual relative matrix]\label{prop:matrix}
For the type-$(1,3)$ system, use the basis of Proposition~\ref{prop:tensor}. On $s_0=1$, with $t=s_1,u=s_2$, the linearized Frobenius on $\Omega^1_{\PP^2}$ has matrix
\begin{equation}\label{eq:matrix}
 \Phi=c\begin{pmatrix}ut^2&1+u^3\\1+t^3&tu^2\end{pmatrix}.
\end{equation}
Its bijective locus is $X^\circ=\PP^2\setminus\{s_0^3+s_1^3+s_2^3=0\}$. The determinant vanishes to order one on this smooth cubic. On the cubic's open $tu\ne0$, the first and stable ranks are one.
\end{proposition}
\begin{proof}
Lemma~\ref{lem:euler-bundle} identifies the relative bundle on all of $X$, including parameters of singular divisors. The Frobenius on ideal cohomology is $h\mapsto s h^2$: the square of $sh$ is $s\cdot(sh^2)$. Serre duality identifies its transpose with the Cartier product map $M_s(v)=\beta(s,v)$. Thus the linearized map on the Euler kernel is the restriction of $M_s^t$ from $(s^{[2]})^\perp$ to $s^\perp$, where $s^{[2]}$ means coefficientwise squaring.

Here
\[
 M_s=c\begin{pmatrix}0&u&t\\u&0&1\\t&1&0\end{pmatrix}.
\]
Use the target frame $(ta+ub,a,b)$ and the actual Frobenius-pullback frame $(t^2a+u^2b,a,b)$. Matrix multiplication gives \eqref{eq:matrix}. Its determinant is $c^2(1+t^3+u^3)$; the other charts give the homogeneous Fermat equation. Its partial derivatives are $s_0^2,s_1^2,s_2^2$, so this cubic is smooth and reduced.

On the cubic, the matrix factors as
\[
                  c\binom{t^2}{u^2}(u\quad t).
\]
The scalar in its next semilinear iterate, apart from a nonzero power of $c$, is
\[
                 ut^4+tu^4=ut(t^3+u^3)=ut.
\]
It is nonzero on the stated open, so that rank-one image is Frobenius bijective.
\end{proof}

\begin{proof}[Proof of Theorem~\ref{thm:three}]
First compute on $X^\circ$. Writing a Lang vector as $c^{-1}$ times a new vector removes $c$ from its fixed-vector equation. Use the matrix without $c$ in \eqref{eq:matrix}.

On the coefficient line $u=0$, the equations are
\[
                x=y^2,\qquad y=(1+t^3)x^2.
\]
A nonzero vector has $y\ne0$, and is parametrized by
\begin{equation}\label{eq:kummer}
                           y^3=(1+t^3)^{-1}.
\end{equation}
This is irreducible over $k(t)$: the right side has valuation $-1$ at each of the three simple zeros of $1+t^3$, whereas the valuation of a cube is divisible by three. Since $k$ contains the cube roots of unity, this is a cyclic degree-three cover. The image on $X^\circ$ therefore contains an element of order three.

On the line $u=t$, the matrix is
\[
             \begin{pmatrix}t^3&1+t^3\\1+t^3&t^3\end{pmatrix},
             \qquad\det=1.
\]
It fixes $(1,1)$. The sum of the two coordinates of a fixed vector satisfies $x+y=(x+y)^2$. The nonzero solution with sum zero is $(1,1)$. For the other two, put $y=x+1$; the first equation becomes
\begin{equation}\label{eq:AS-slice}
                            x^2+x=1+t^3.
\end{equation}
No rational function $h$ can satisfy this equation, because every pole of $h^2+h$ has even order while the right side has a pole of order three at infinity. Thus this restriction contains an element of order two, interchanging the other two nonzero vectors. Both restrictions are actual restrictions of one local system; their images lie in its image up to conjugacy.

A subgroup of $S_3=\GL_2(\FF_2)$ containing elements of orders two and three is all of $S_3$. This proves the full image on $X^\circ$. Lemma~\ref{lem:smooth} gives a nonempty smooth-member open, and it meets $X^\circ$ because the parameter plane is integral. Restriction to a nonempty open of a normal integral base preserves a finite monodromy image: each connected finite \'etale covering component is normal and integral, and stays integral on the smaller open. The calculation therefore applies to the smooth curve family, even though the two coefficient lines themselves need not consist of smooth divisors.

Transitivity on the three nonzero vectors proves connectedness of the nonzero cover. Lemma~\ref{lem:euler-bundle} gives genus four and two-rank $f(A)+2$.
\end{proof}

\section{The explicit primitive lines and the branch module}\label{sec:double}

We prepare the geometric family of Theorem~\ref{thm:odd} and prove the natural cohomology representation of a wild double cover. This representation, rather than a numerical rank formula alone, will determine monodromy.

\begin{lemma}[The lines $L_m$]\label{lem:explicit-lines}
For every $m\ge1$, the line in \eqref{eq:lines} is ample and globally generated of type $(1,2m+1)$, with the numerical invariants in Theorem~\ref{thm:odd}. The surface $E^2$ has two-rank zero.
\end{lemma}
\begin{proof}
The projective equation of $E$ is $Y^2Z+YZ^2=X^3$. Its partials are $X^2,Z^2,Y^2$, with no common projective zero. Its invariant differential $dx$ is exact in the function field and is killed by Cartier. By duality, Frobenius on $H^1(E,\OO_E)$ is zero. Thus $f(E)=0$ and $f(E^2)=0$.

The elliptic divisors $V,H,\Gamma$ have self-intersection zero, and each pair has intersection one. They are fibres of the two projections and addition; the two defining homomorphisms at each intersection have integral determinant $\pm1$. Therefore
\[
                         L_m^2=4m+2,\qquad L_m\cdot V=2.
\]
We prove global generation directly. For $v\in E(k)$ consider
\[
 D_v=m(\{v\}\times E)+(E\times\{mv\})
                      +\{(a,b):a+b=-mv\}.
\]
Its line is $L_m$: in $\Pic^0(E)$, the three differences from the divisors in \eqref{eq:lines} contribute $mv$ on the first factor, $mv$ on the second, and $-mv$ on each factor. These cancel by the addition identity for degree-zero lines \citep[Proposition~I.8.4 and Remark~I.8.5, pp.~35--36]{Milne}. Given $(a,b)$, the divisor contains it only if $v=a$, $mv=b$, or $-mv=a+b$. Each excludes finitely many values because $[m]$ is a finite isogeny, including when $m$ is even. Some $v$ avoids all three, so $L_m$ is globally generated.

For every integral curve $T\subset E^2$, the intersections with $V$ and $H$ are the nonnegative degrees of the two projections, with degree zero for a constant projection. They cannot both be zero. The intersection with $\Gamma$ is likewise nonnegative. Hence $L_m\cdot T>0$. The morphism defined by the globally generated line contracts no curve. Every positive-dimensional projective fibre contains a curve, so the morphism is finite. Its pullback of $\OO(1)$ is ample, proving ampleness of $L_m$. Vanishing and Riemann--Roch give $h^0(L_m)=2m+1$.

Relative to the product degree-one principal polarization, the polarization homomorphism is represented by
\begin{equation}\label{eq:polarization-matrix}
                         \begin{pmatrix}m+1&1\\1&2\end{pmatrix}.
\end{equation}
The summands $mV,H$ give the diagonal part, and the addition pullback gives the all-one matrix. Its determinant is $2m+1$, and its entries have greatest common divisor one, so its Smith factors are $1,2m+1$. Integral unimodular row and column operations act by automorphisms of $E^2$. Its kernel is therefore $E[2m+1]$, giving type $(1,2m+1)$. The polarization identities used in this matrix calculation follow from the theorem of the square and the addition formula \citep[Theorem~I.5.5 and Proposition~I.8.4, pp.~23 and 35--36]{Milne}.
\end{proof}

\begin{lemma}[The natural branch permutation module]\label{lem:branch}
Let $C$ be smooth, proper and connected, and let $f:C\to E$ be a finite separable morphism of degree two. If $B$ is its reduced geometric branch set, then $B\ne\varnothing$ and there is a natural isomorphism
\begin{equation}\label{eq:branch-module}
                    H^1_{\et}(C,\FF_2)\simeq\FF_2^B/\FF_2\mathbf1.
\end{equation}
It commutes with field automorphisms and with isomorphisms of covers. In particular $f(C)=|B|-1$.
\end{lemma}
\begin{proof}
Artin--Schreier gives $H^1_{\et}(E,\FF_2)=0$ and $H^2_{\et}(E,\FF_2)=0$: Frobenius on $H^1(E,\OO_E)$ is zero, so $F-1=-1$ is bijective, and coherent $H^2$ of the curve vanishes. If $B$ were empty, the connected degree-two \'etale cover would define a nonzero class in the first group. Thus $B$ is nonempty.

A separable quadratic extension in characteristic two is Galois. Let $\iota$ be its unique nontrivial automorphism and $j:E\setminus B\hookrightarrow E$. There is an exact sequence of \'etale sheaves
\begin{equation}\label{eq:branch-sheaf}
 0\longrightarrow\FF_{2,E}\longrightarrow f_*\FF_{2,C}
          \xrightarrow{1+\iota}j_!\FF_{2,E\setminus B}\longrightarrow0.
\end{equation}
Off $B$, the first map is the diagonal into $\FF_2^2$ and the last map is addition. Over a branch point the geometric fibre has one point, pullback is an isomorphism and addition is zero. The image of $1+\iota$ lies in the invariant constant subsheaf and is exactly the indicated extension by zero. This proves exactness on geometric stalks, including wild ramification.

For a finite morphism, higher direct images of an abelian \'etale sheaf vanish, and its direct-image stalk is the direct sum over the finite geometric fibre \citep[Lemma~59.55.1 and Proposition~59.55.2, full proofs]{StacksFinite}. These statements impose no invertibility condition on the coefficient prime. Their proof decomposes the finite scheme over a strict henselization into strictly henselian local schemes, where sections equal the closed-point stalk. Hence Leray and \eqref{eq:branch-sheaf} give
\[
 H^1_{\et}(C,\FF_2)\simeq H^1_{\et}(E,j_!\FF_2).
\]
The exact sequence $0\to j_!\FF_2\to\FF_{2,E}\to i_*\FF_{2,B}\to0$, where $i:B\hookrightarrow E$, identifies the right side with the cokernel of diagonal constants $\FF_2\to\FF_2^B$. All maps are canonical for the cover, and its involution is unique. This proves the claimed naturality.
\end{proof}

\section{The complete elliptic branch series}\label{sec:series}

We identify the varying branch labels for every odd $n$, without a finite-prime generic-rank assumption. The decisive fact is that the entire coefficient system descends to a complete linear series on an elliptic curve after relative Frobenius.

\begin{proposition}[Frobenius descent of the branch series]\label{prop:series}
Fix $m\ge1$, put $n=2m+1$ and $W=H^0(A,L_m)$. There is a degree-$n$ line $N$ on an elliptic curve $R^{(2)}$ and a projective linear identification
\begin{equation}\label{eq:series-isom}
                             |L_m|\simeq|N|
\end{equation}
with the following property. If a smooth member $C_s$ corresponds to a divisor $D_s\in|N|$, its degree-two map to the first factor $E$ has branch points naturally identified with the geometric support of $D_s$. The identification is equivariant for coefficient-field automorphisms. Consequently
\begin{equation}\label{eq:every-rank}
                       f(C_s)=|\operatorname{Supp}D_s|-1.
\end{equation}
The finite \'etale permutation system on the open of reduced divisors in $|N|$ supplies these branch labels.
\end{proposition}
\begin{proof}
Let $q:A\to E$ be the first projection and $\varphi:A\to\PP(W^*)$ the complete-system morphism. On $q^{-1}(a)$ the line restricts to $\OO_E(O+(-a))$, of degree two. The restriction of $W$ is its full two-dimensional section space: it is basepoint free by global generation, and a one-dimensional space cannot generate a positive-degree line. Thus $\varphi$ on each fibre is its complete degree-two elliptic map followed by a linear embedding of $\PP^1$.

That map is separable. A purely inseparable degree-two map from an elliptic curve to $\PP^1$ would factor through relative Frobenius and a degree-one map, contradicting genus; see \citet[Lemmas~53.13.3--53.13.5, with proofs]{StacksFrob}. Its involution is $b\mapsto-a-b$, since a divisor in $|O+(-a)|$ has elliptic group sum $-a$. Its ramification support consists exactly of the fixed points $a+2b=0$: away from them the two points in a fibre are distinct, whereas a fixed point is a double point of that fibre. This is valid in wild characteristic two; on smooth curves ramification is detected by the differential \citep[Lemma~53.12.1]{StacksRam}.

The map $[2]:E\to E$ is radicial of degree four because $E$ is supersingular. The reduced fixed-point curve is the smooth graph
\begin{equation}\label{eq:R}
                    i:R=E\hookrightarrow A,\qquad
                    b\longmapsto(-[2]b,b).
\end{equation}
Its tangent is vertical for $q$, since $d[2]=0$. The differential of the complete fibre map kills that tangent at every point of $R$. Thus $d(\varphi\circ i)=0$.

We spell out the resulting scheme-theoretic factorization. For a smooth curve over perfect $k$, the kernel of $d:k(R)\to\Omega^1_{k(R)/k}$ is $k(R)^2$. Hence the coordinate ratios of a zero-differential morphism to projective space belong to the relative-Frobenius subfield. Their roots are regular where required, by the valuations in the smooth curve's local discrete valuation rings. They define a morphism on $R^{(2)}$; the unique local factorizations glue. Equivalently, extend the generic factorization to the smooth proper curve by properness of projective space. This is precisely the factorization mechanism in \citet[Lemma~53.13.8, full proof]{StacksFrob}. We obtain
\[
 R\xrightarrow{F_R}R^{(2)}\xrightarrow{\psi}\PP(W^*),
 \qquad \varphi\circ i=\psi\circ F_R.
\]
Set $N=\psi^*\OO(1)$. Then $F_R^*N\simeq i^*L_m$. The degrees of the restrictions of $V,H,\Gamma$ to \eqref{eq:R} are $4,1,1$, respectively. Therefore
\begin{equation}\label{eq:restriction-degree}
                         \deg i^*L_m=4m+2=2n,
                         \qquad\deg N=n.
\end{equation}

Restriction $W\to H^0(R,i^*L_m)$ is injective. Indeed
\[
 (L_m-R)\cdot L_m=L_m^2-L_m\cdot R=2n-2n=0.
\]
A nonzero effective divisor has positive intersection with the ample $L_m$. The line $L_m(-R)$ is not trivial, since $R$, a fibre of the smooth homomorphism $(a,b)\mapsto a+2b$, has $R^2=0$, and $(L_m-R)^2=-2n$. It therefore has no nonzero section. The $n$ projective coordinate sections of $N$ pull back to the independent restrictions of $W$, so they are independent. Elliptic Riemann--Roch gives $h^0(N)=n$; they consequently form the complete section space. This proves the actual linear-series identification \eqref{eq:series-isom}, not just equality of its dimensions. For corresponding sections,
\begin{equation}\label{eq:restriction-divisor}
                             C_s|_R=F_R^*D_s.
\end{equation}

A smooth connected $C_s$ maps finitely to $E$ under $q$, with degree $L_m\cdot V=2$. It cannot be a fibre, whose genus would be one. Its genus is $n+1$ by Lemma~\ref{lem:euler-bundle}; a purely inseparable degree-two map would preserve genus, so the map is separable.

Its ramification support is exactly $C_s\cap R$. At a point of $C_s$, ramification means that the tangent of $C_s$ is vertical, or equivalently that its defining equation restricted to the vertical fibre has a multiple zero. That restricted section is a nonzero degree-two section, because a smooth connected ample member cannot contain a vertical fibre. Its zeros coincide exactly at a fixed point of $b\mapsto-a-b$, namely on $R$. There are no other ramification points.

The radicial map $q\circ i=[-2]$ factors as
\[
                         [-2]=h\circ F_R,
                         \qquad h:R^{(2)}\to E,
\]
where $h$ is a radicial isogeny of degree two. Its existence follows from the same zero-differential factorization; its degree and radiciality follow from those of $[-2]$ and $F_R$. Equations \eqref{eq:restriction-divisor} and this factorization identify the branch set with $h(\operatorname{Supp}D_s)$. A radicial map is universally injective and induces a bijection on geometric points onto its image. This identification commutes with coefficient-field automorphisms.

Let $V_N\subset|N|$ be the open of reduced divisors. It is nonempty: choose $n-1$ general distinct elliptic points and take the last one to make their sum equal the class of $N$; coincidences impose proper closed conditions. The universal divisor over $V_N$ is finite \'etale of degree $n$. Its geometric points provide the branch labels on the corresponding smooth-member open. No assertion that the scheme-theoretic image under the inseparable map $h$ is itself finite \'etale is needed. The label system is the universal reduced divisor, and its geometric-point identification is the one used in Lemma~\ref{lem:branch}. That lemma proves \eqref{eq:every-rank}.
\end{proof}

\section{Exact permutation monodromy and Lang components}\label{sec:odd-proof}

We finish the infinite-family theorem by computing the monodromy of a complete elliptic linear series. The calculation gives transitivity inside each component, as well as the component count.

\begin{lemma}[Ordered divisors in a complete elliptic series]\label{lem:ordered}
Let $N$ be a line of degree $n\ge3$ on an elliptic curve $B/k$. Over the open $V_N\subset|N|$ of reduced divisors, the permutation image on the universal degree-$n$ divisor is $S_n$.
\end{lemma}
\begin{proof}
Identify $\Pic^n(B)$ with $B$ using an origin. A divisor $x_1+\cdots+x_n$ has class $N$ exactly when the group sum of its points is a specified point $a_N$. The ordered tuples with that sum form a scheme isomorphic to $B^{n-1}$: the final coordinate is $a_N-\sum_{i<n}x_i$. Removing all pairwise diagonals leaves a nonempty integral open $T$.

The map $T\to V_N$ is the cover ordering the distinct points of the universal divisor. It is a finite \'etale $S_n$-torsor. This remains true when the characteristic divides $n!$: the symmetric group here is the constant \'etale group and its action on orderings is free. The total space $T$ is connected, so the image of the fundamental group is all of $S_n$.
\end{proof}

\begin{proof}[Proof of Theorem~\ref{thm:odd}]
Lemma~\ref{lem:explicit-lines} gives the numerical assertions. Lemma~\ref{lem:smooth} supplies smooth members for every odd $n$. Under Proposition~\ref{prop:series}'s projective linear identification, this nonempty smooth open meets $V_N$. On their intersection the branch set has exactly $n$ points, so Lemma~\ref{lem:branch} gives $f(C)=n-1$. For every smooth member, \eqref{eq:every-rank} gives $f(C)\le n-1$, with equality exactly when $D_s$ is reduced. Thus this intersection is precisely the smooth maximal-rank open $U_m$. As $f(A)=0$, Lemma~\ref{lem:euler-bundle} identifies its full \'etale system with the relative Lang system.

Lemma~\ref{lem:ordered} gives the full permutation image $S_n$ on $V_N$, and restriction to the nonempty open $U_m$ preserves it. The natural branch-module isomorphism identifies the resulting representation with \eqref{eq:Sn}. For clarity, this can be checked at the geometric generic point: the branch labels and the reduced-divisor labels are the same Galois set under the radicial map $h$, and the generic Galois image is the image of a finite local system on a normal integral parameter space. This avoids any replacement of an inseparable branch-support scheme by an \'etale one.

The permutation action on $\FF_2^n/\FF_2\mathbf1$ is faithful for odd $n\ge3$. If a permutation fixes the class of every basis vector, its difference from that basis vector is either zero or $\mathbf1$. A nonzero difference of two distinct basis vectors has weight two, which is not $n$. Therefore every basis vector is fixed.

A quotient class is a subset of $\{1,\ldots,n\}$ up to complementation. Because $n$ is odd, every nonzero class has a unique representative of size $j$ with $1\le j\le(n-1)/2$. The full $S_n$ acts transitively on each size class. These are exactly its nonzero orbits, with cardinalities \eqref{eq:degrees}. They give exactly the claimed connected components. For $n=3$ there is one orbit and $S_3=\GL_2(\FF_2)$. For $n\ge5$ there are at least two, whereas the full linear group is transitive on all nonzero vectors. Hence the subgroup is proper.
\end{proof}

\paragraph{A small case by hand.}
For $n=5$, the nonzero quotient vectors are singleton classes and two-element classes. Complementation has already identified weights four with one and three with two. Their degrees are $5$ and $10$, adding to $2^4-1=15$. The faithful image has order $5!=120$; the full linear group has order
\[
                      (16-1)(16-2)(16-4)(16-8)=20160.
\]
For $n=3$, singleton classes account for all $2^2-1=3$ nonzero vectors, explaining why the product example agrees with Theorem~\ref{thm:three}.

\section{Scope, exact calculations, and method}\label{sec:scope}

We specify what these image calculations decide and what they leave open. Theorem~\ref{thm:three} holds for every fixed nonordinary type-$(1,3)$ surface under global generation. Its proof includes the geometric nonvanishing of the Cartier tensor and the existence of smooth members. Theorem~\ref{thm:odd} treats every odd degree in the explicit product family; it does not depend on testing a finite list of primes. The branch-series construction proves both the rank and the exact image there.

These statements leave the full relative image for a general fixed nonordinary abelian surface and an arbitrary primitive polarization undecided. In particular, an all-prime generic dual-Cartier-rank statement for unrelated primitive systems is not supplied by the explicit family. No conclusion about a full image in every sufficiently large degree of a fixed projective embedding follows. The failure of very ampleness in \eqref{eq:lines} is concrete: its restriction to $V$ is a degree-two elliptic line and cannot embed that fibre.

The \emph{Primitive monodromy certificate}, available through the workbench, checks the displayed Euler determinant and both Lang restrictions by sparse polynomial arithmetic over $\FF_2$. A separate cycle-index calculation gives the nonzero permutation-orbit counts for $n=3,5,7,11,13,17$, and a generated-matrix calculation gives the faithful $120$-element $S_5$ image and its orbits of sizes five and ten. Earlier exact programs were read and rerun, including their full Euler-frame, wall-iterate, Smith-matrix and permutation checks. These calculations corroborate finite algebra; the Frobenius descent, completeness, infinite-family statement and monodromy identifications are proved in the text.

\paragraph{Method and availability.}
The mathematics and text were produced by GPT-6 Astra Pro. Full author mathematical and expository readings, individual page inspection, and a separate full proof and source comparison are complete. A separate same-model reconstruction and comparison established the fixed type-$(1,3)$ image and the image bounds for $n=5,7,11,13,17$, relative to explicitly admitted geometric and rank inputs. A second complete reconstruction and comparison establishes the all-odd exact-image theorem, using invariant rational sections, a complete square-root series and a separate first-jet smoothness argument. Exact programs were checked separately. These are same-model checks, not formal verification; no human has checked the mathematics. The new text and code are dedicated under CC0~1.0. Editable sources, build instructions and calculations are available through the \href{\workbenchurl}{schemes-and-moduli workbench}; the \href{\recordurl}{mathematics-workbenches index} records versions and their provenance.

\begin{thebibliography}{99}
\bibitem[Catanese and Sernesi(2022)]{CataneseSernesi}
Catanese, F., \& Sernesi, E. (2022, December 30). \emph{The Hesse pencil and polarizations of type $(1,3)$ on Abelian surfaces} (Version 1). arXiv:2212.14877. \url{https://arxiv.org/abs/2212.14877v1}.
\bibitem[Edixhoven et~al.(n.d.)]{EGM}
Edixhoven, B., van der Geer, G., \& Moonen, B. (n.d.). \emph{Abelian varieties: Preliminary version of the first chapters}. Author manuscript. \url{https://van-der-geer.nl/~gerard/AV.pdf}.
\bibitem[Grothendieck(n.d.)]{Grothendieck}
Grothendieck, A. (n.d.). \emph{Notes on crystalline cohomology}, cote~4, \S8.6, archive sheet~11, PDF image~12. Institut Montpelli\'erain Alexander Grothendieck. \url{https://grothendieck.umontpellier.fr/4.pdf}.
\bibitem[Illusie(1979)]{Illusie}
Illusie, L. (1979). Complexe de de Rham--Witt et cohomologie cristalline. \emph{Annales scientifiques de l'\'Ecole normale sup\'erieure, 12}(4), 501--661. \url{https://doi.org/10.24033/asens.1374}.
\bibitem[Milne(2008)]{Milne}
Milne, J. S. (2008, March 16). \emph{Abelian varieties} (Version 2.00). Course notes. \url{https://www.jmilne.org/math/CourseNotes/AVc.pdf}.
\bibitem[The Stacks Project Authors(2026a)]{StacksAS}
The Stacks Project Authors. (2026a). \emph{The Artin--Schreier sequence}. Section~59.63, especially Lemma~59.63.2, tag~0A3J. \url{https://stacks.math.columbia.edu/tag/0A3J}.
\bibitem[The Stacks Project Authors(2026b)]{StacksFinite}
The Stacks Project Authors. (2026b). \emph{Vanishing of finite higher direct images}. Lemma~59.55.1 and Proposition~59.55.2, tag~03QN. \url{https://stacks.math.columbia.edu/tag/03QN}.
\bibitem[The Stacks Project Authors(2026c)]{StacksFrob}
The Stacks Project Authors. (2026c). \emph{Inseparable maps}. Lemmas~53.13.3--53.13.5 and 53.13.8, tags~0CCV and 0CD3. \url{https://stacks.math.columbia.edu/tag/0CCV}; \url{https://stacks.math.columbia.edu/tag/0CD3}.
\bibitem[The Stacks Project Authors(2026d)]{StacksRam}
The Stacks Project Authors. (2026d). \emph{Riemann--Hurwitz}, Lemma~53.12.1, tag~0C1C. \url{https://stacks.math.columbia.edu/tag/0C1C}.
\bibitem[Sun(2006)]{Sun}
Sun, X. (2006, August 18). \emph{Stability of direct images under Frobenius morphism} (Version 2). arXiv:math/0608043. \url{https://arxiv.org/abs/math/0608043v2}.
\bibitem[Zou(2026)]{Zou}
Zou, H. (2026, September 29). \emph{Quasi-$F$-split primitive symplectic varieties in positive characteristic}. Author preprint. \url{https://www.math.uni-bielefeld.de/~hzou/papers/qfpsv/}.
\end{thebibliography}
\end{document}
